% Options for packages loaded elsewhere
\PassOptionsToPackage{unicode}{hyperref}
\PassOptionsToPackage{hyphens}{url}
\PassOptionsToPackage{dvipsnames,svgnames,x11names}{xcolor}
\documentclass[
  10pt,
  fontset=fandol]{ctexart}
\usepackage{xcolor}
\usepackage[margin=16mm]{geometry}
\usepackage{amsmath,amssymb}
\setcounter{secnumdepth}{-\maxdimen} % remove section numbering
\usepackage{iftex}
\ifPDFTeX
  \usepackage[T1]{fontenc}
  \usepackage[utf8]{inputenc}
  \usepackage{textcomp} % provide euro and other symbols
\else % if luatex or xetex
  \usepackage{unicode-math} % this also loads fontspec
  \defaultfontfeatures{Scale=MatchLowercase}
  \defaultfontfeatures[\rmfamily]{Ligatures=TeX,Scale=1}
\fi
\usepackage{lmodern}
\ifPDFTeX\else
  % xetex/luatex font selection
\fi
% Use upquote if available, for straight quotes in verbatim environments
\IfFileExists{upquote.sty}{\usepackage{upquote}}{}
\IfFileExists{microtype.sty}{% use microtype if available
  \usepackage[]{microtype}
  \UseMicrotypeSet[protrusion]{basicmath} % disable protrusion for tt fonts
}{}
\makeatletter
\@ifundefined{KOMAClassName}{% if non-KOMA class
  \IfFileExists{parskip.sty}{%
    \usepackage{parskip}
  }{% else
    \setlength{\parindent}{0pt}
    \setlength{\parskip}{6pt plus 2pt minus 1pt}}
}{% if KOMA class
  \KOMAoptions{parskip=half}}
\makeatother
\setlength{\emergencystretch}{3em} % prevent overfull lines
\providecommand{\tightlist}{%
  \setlength{\itemsep}{0pt}\setlength{\parskip}{0pt}}
\usepackage{amsmath,amssymb,mathtools,needspace}
\xeCJKDeclareCharClass{CJK}{"2032,"0400->"04FF,"2160->"216B,"2460->"2473,"25A0,"25C7,"25CB,"25CF}
\IfFontExistsTF{Arial Unicode MS}{\xeCJKsetup{AutoFallBack=true}\setCJKfallbackfamilyfont{\CJKrmdefault}{Arial Unicode MS}}{}
\setcounter{secnumdepth}{0}
\setlength{\emergencystretch}{2em}
\allowdisplaybreaks
\usepackage{bookmark}
\IfFileExists{xurl.sty}{\usepackage{xurl}}{} % add URL line breaks if available
\urlstyle{same}
\hypersetup{
  pdftitle={插值余项},
  colorlinks=true,
  linkcolor={Maroon},
  filecolor={Maroon},
  citecolor={Blue},
  urlcolor={Blue},
  pdfcreator={LaTeX via pandoc}}

\title{插值余项}
\author{}
\date{}

\begin{document}
\maketitle

\subsubsection{2.2.4 插值余项}\label{ux63d2ux503cux4f59ux9879}

若在 \([a,b]\) 上用 \(L_n(x)\) 近似 \(f(x)\)，则其截断误差为 \(R_n(x)=f(x)-L_n(x)\)，\(R_n(x)\) 也称为插值多项式的余项或插值余项。关于插值余项估计有以下定理。

\textbf{定理 2.2} 设 \(f^{(n)}(x)\) 在 \([a,b]\) 上连续，\(f^{(n+1)}(x)\) 在 \((a,b)\) 内存在，节点 \(a\le x_0<x_1<\cdots<x_n\le b\)，\(L_n(x)\) 是满足条件式（2.2.8）的插值多项式，则对于任何 \(x\in[a,b]\)，插值余项

\[
R_n(x)=f(x)-L_n(x)=\frac{f^{(n+1)}(\xi)}{(n+1)!}\omega_{n+1}(x).
\tag{2.2.14}
\]

这里，\(\xi\in(a,b)\) 且依赖于 \(x\)，\(\omega_{n+1}(x)\) 是式（2.2.12）所定义的。

\textbf{证明} 由给定条件知，\(R_n(x)\) 在节点 \(x_k\)（\(k=0,1,\cdots,n\)）上为零，即

\[
R_n(x_k)=0\qquad(k=0,1,\cdots,n),
\]

于是

\[
R_n(x)=K(x)(x-x_0)(x-x_1)\cdots(x-x_n)=K(x)\omega_{n+1}(x),
\tag{2.2.15}
\]

其中 \(K(x)\) 是与 \(x\) 有关的待定函数。

现把 \(x\) 看成 \([a,b]\) 上一个固定点，作函数

\[
\varphi(t)=f(t)-L_n(t)-K(x)(t-x_0)(t-x_1)\cdots(t-x_n),
\]

根据插值条件及余项定义，可知 \(\varphi(t)\) 在点 \(x_0,x_1,\cdots,x_n\) 及 \(x\) 处均为零，故 \(\varphi(t)\) 在 \([a,b]\) 上有 \(n+2\) 个零点，根据 Rolle 定理，\(\varphi'(t)\) 在 \(\varphi(t)\) 的两个零点间至少有一个零点，故 \(\varphi'(t)\) 在 \([a,b]\) 上至少有 \(n+1\) 个零点。对 \(\varphi'(t)\) 再应用 Rolle 定理，可知 \(\varphi''(t)\) 在 \([a,b]\) 上至少有 \(n\) 个零点。依此类推，\(\varphi^{(n+1)}(t)\) 在 \((a,b)\) 内至少有一个零点，记作 \(\xi\in(a,b)\)，使

\[
\varphi^{(n+1)}(\xi)=f^{(n+1)}(\xi)-(n+1)!K(x)=0,
\]

于是

\[
K(x)=\frac{f^{(n+1)}(\xi)}{(n+1)!},\qquad \xi\in(a,b)
\]

且依赖于 \(x\)。

将上式代入式（2.2.15），就得到余项表达式（2.2.14）。证毕。

应当指出，余项表达式只有在 \(f(x)\) 的高阶导数存在时才能应用。\(\xi\) 在 \((a,b)\) 内的具体位置通常不可能给出，如果可以求出 \(\max_{a\le x\le b}|f^{(n+1)}(x)|=M_{n+1}\)，那么插值多项式 \(L_n(x)\) 逼近 \(f(x)\) 的截断误差限是

\[
|R_n(x)|\le\frac{M_{n+1}}{(n+1)!}|\omega_{n+1}(x)|.
\tag{2.2.16}
\]

当 \(n=1\) 时，线性插值余项为

\href{../../source-images/pdf-033.jpeg}{原页：书页 20／PDF 33}

\begin{quote}
页首续 PDF 32 的线性插值余项；末行误差估计的不等式续 PDF 34。
\end{quote}

\[
R_1(x)=\frac12 f''(\xi)\omega_2(x)
=\frac12 f''(\xi)(x-x_0)(x-x_1),
\qquad \xi\in[x_0,x_1];
\tag{2.2.17}
\]

当 \(n=2\) 时，抛物插值的余项为

\[
R_2(x)=\frac16 f'''(\xi)(x-x_0)(x-x_1)(x-x_2),
\qquad \xi\in[x_0,x_2].
\tag{2.2.18}
\]

\textbf{例 2.1} 已知 \(\sin0.32=0.314\,567\)，\(\sin0.34=0.333\,487\)，\(\sin0.36=0.352\,274\)，用线性插值及抛物插值计算 \(\sin0.336\,7\) 的值，并估计截断误差。

\textbf{解} 由题意取 \(x_0=0.32\)，\(y_0=0.314\,567\)，\(x_1=0.34\)，\(y_1=0.333\,487\)，\(x_2=0.36\)，\(y_2=0.352\,274\)。

用线性插值计算，取 \(x_0=0.32\) 及 \(x_1=0.34\)，由式（2.2.3）得

\[
\begin{aligned}
\sin0.336\,7\approx L_1(0.336\,7)
&=y_0+\frac{y_1-y_0}{x_1-x_0}(0.336\,7-x_0)\\
&=0.314\,567+\frac{0.018\,92}{0.02}\times0.016\,7\\
&=0.330\,365.
\end{aligned}
\]

其截断误差由式（2.2.17）得

\[
|R_1(x)|\le\frac{M_2}{2}|(x-x_0)(x-x_1)|,
\]

其中

\[
M_2=\max_{x_0\le x\le x_1}|f''(x)|.
\]

因

\[
f(x)=\sin x,\qquad f''(x)=-\sin x,
\]

可取

\[
M_2=\max_{x_0\le x\le x_1}|\sin x|=\sin x_1\le0.333\,5,
\]

于是

\[
\begin{aligned}
|R_1(0.336\,7)|
&=|\sin0.336\,7-L_1(0.336\,7)|\\
&\le\frac12\times0.333\,5\times0.016\,7\times0.003\,3\\
&\le0.92\times10^{-5}.
\end{aligned}
\]

若取 \(x_1=0.34\)，\(x_2=0.36\) 为节点，则线性插值为

\[
\begin{aligned}
\sin0.336\,7\approx\widetilde L_1(0.336\,7)
&=y_1+\frac{y_2-y_1}{x_2-x_1}(0.336\,7-x_1)\\
&=0.333\,487+\frac{0.018\,787}{0.02}\times(-0.003\,3)\\
&=0.330\,387,
\end{aligned}
\]

其截断误差为

\[
|\widetilde R_1(x)|\le\frac{M_2}{2}|(x-x_1)(x-x_2)|,
\]

其中

\[
M_2=\max_{x_0\le x\le x_2}|f''(x)|\le0.352\,3.
\]

于是

\[
|\widetilde R_1(0.336\,7)|=|\sin0.336\,7-\widetilde L_1(0.336\,7)|
\]

\href{../../source-images/pdf-034.jpeg}{原页：书页 21／PDF 34}

\begin{quote}
页首续 PDF 33 的 \(|\widetilde R_1(0.3367)|\) 误差估计。
\end{quote}

\[
\begin{aligned}
&\le\frac12\times0.352\,3\times0.003\,3\times0.023\,3\\
&\le1.36\times10^{-5}.
\end{aligned}
\]

用抛物插值计算 \(\sin0.336\,7\) 时，由式（2.2.7）得

\[
\begin{aligned}
\sin0.336\,7\approx{}&
y_0\frac{(x-x_1)(x-x_2)}{(x_0-x_1)(x_0-x_2)}
+y_1\frac{(x-x_0)(x-x_2)}{(x_1-x_0)(x_1-x_2)}\\
&+y_2\frac{(x-x_0)(x-x_1)}{(x_2-x_0)(x_2-x_1)}\\
={}&L_2(0.336\,7)\\
={}&0.314\,567\times\frac{0.768\,9\times10^{-4}}{0.000\,8}
+0.333\,487\times\frac{3.89\times10^{-4}}{0.000\,4}\\
&+0.352\,274\times\frac{-0.551\,1\times10^{-4}}{0.000\,8}\\
={}&0.330\,374.
\end{aligned}
\]

这个结果与六位有效数字的正弦函数表完全一样，这说明查表时用二次插值精度已相当高了。其截断误差限由式（2.2.18）得

\[
|R_2(x)|\le\frac{M_3}{6}|(x-x_0)(x-x_1)(x-x_2)|,
\]

其中

\[
M_3=\max_{x_0\le x\le x_2}|f'''(x)|=\cos x_0<0.828,
\]

于是

\[
\begin{aligned}
|R_2(0.336\,7)|
&=|\sin0.336\,7-L_2(0.336\,7)|\\
&\le\frac16\times0.828\times0.016\,7\times0.003\,3\times0.023\,3\\
&\le0.178\times10^{-7}.
\end{aligned}
\]

\begin{center}\rule{0.5\linewidth}{0.5pt}\end{center}

\subsection{相关原页校注（agent 补充）}\label{ux76f8ux5173ux539fux9875ux6821ux6ce8agent-ux8865ux5145}

以下校注来自本节覆盖的原页，可能也涉及同页相邻小节；原文照录与校正说明分开保留。

\subsubsection{原书 PDF 页 34 的校注}\label{ux539fux4e66-pdf-ux9875-34-ux7684ux6821ux6ce8}

\href{../pages/pdf-034.md}{完整原页转录} · \href{../../source-images/pdf-034.jpeg}{原页图像}

校注记录：NA5E-E002。

\subsubsection{校注（agent 补充）：NA5E-E002}\label{ux6821ux6ce8agent-ux8865ux5145na5e-e002}

上面的 \(0.828\)、\(0.178\times10^{-7}\) 和代入行的 \(3.89\times10^{-4}\) 均经过原页放大核对，确为扫描原书印字。\textbf{这几行不能直接作为正确的数值结果使用。}

\begin{enumerate}
\def\labelenumi{\arabic{enumi}.}
\tightlist
\item
  本例以弧度计算，\(\cos0.32\approx0.9492354180824409\)，所以 \(\cos x_0<0.828\) 不成立。可用保守上界 \(M_3<0.95\)。
\item
  即使照原书使用 \(0.828\)，该乘积也等于 \(1.77200694\times10^{-7}\)，大于原文所写的 \(0.178\times10^{-7}=1.78\times10^{-8}\)。
\item
  对正确的导数上界，用原节点得到的截断误差上界为
\end{enumerate}

\[
|R_2(0.3367)|
\le\frac{0.95}{6}\,(0.0167)(0.0033)(0.0233)
=2.03309975\times10^{-7}.
\]

这里的 \(R_2\) 指用精确节点函数值定义的插值截断误差。

\begin{enumerate}
\def\labelenumi{\arabic{enumi}.}
\setcounter{enumi}{3}
\tightlist
\item
  二次插值代入行的中间乘积应为 \(0.0167\times0.0233=0.00038911=3.8911\times10^{-4}\)。若严格按原书印出的 \(3.89\times10^{-4}\) 计算整行，会得到 \(0.3302826531125\)，不能得到末行的 \(0.330374\)。保留精确乘积，并使用题目给的有限精度表值，可得
\end{enumerate}

\[
\begin{aligned}
L_1(0.3367)&=0.3303652,\\
\widetilde L_1(0.3367)&=0.330387145,\\
L_2(0.3367)&=0.3303743620375.
\end{aligned}
\]

最后一项保留六位有效数字确实为 \(0.330374\)，与 \(\sin0.3367\approx0.3303741915556280\) 的六位有效数字一致。

\subsection{本节覆盖原页的校注记录}\label{ux672cux8282ux8986ux76d6ux539fux9875ux7684ux6821ux6ce8ux8bb0ux5f55}

下列记录按原页关联，含同页相邻小节的校注；计算时同时读取上文校注和完整原页。

\begin{itemize}
\tightlist
\item
  \texttt{NA5E-E002}：原书 PDF 页 34
\end{itemize}

\end{document}
